% Affine registration -- version written with Claude.
% Pulled in by sections/registration/affine-registration.tex.

\subsubsection{Claude}           % first version of the affine-registration text
\label{sec:affine_claude}        % \secref{sec:affine_claude}

% --- Values to fill in (shown in red in the PDF until replaced) ---
\newcommand{\ncases}{\TODO{N}}                % number of AIIB23 segmentations registered
\newcommand{\antspyversion}{\TODO{x.y.z}}     % python -c "import ants; print(ants.__version__)"

% ------------------------------------------------------------
\paragraph{Objective.}           % why the registration is performed

Airway trees extracted from different subjects differ in position, orientation
and overall size, as a consequence of patient positioning in the scanner,
differences in body size and acquisition parameters. To make them directly
comparable, all airway segmentations were mapped into a common spatial
reference frame defined by a template airway tree. An affine transformation was
chosen for this purpose because it compensates for these global differences
while preserving the local geometry of each airway tree, which is the object of
the subsequent analysis \cite{maintz1998survey}.

% ------------------------------------------------------------
\paragraph{Data.}                % input images and reference template

The input data consisted of the \ncases{} ground-truth airway segmentations of
the AIIB23 challenge dataset \cite{nan2024aiib23}, stored as binary
three-dimensional volumes in NIfTI format. A single airway tree (template~22)
was used as the reference image for all registrations.
% TODO: add one sentence on how template 22 was selected or constructed.

% ------------------------------------------------------------
\paragraph{Transformation model.}  % mathematical definition of the affine map

Registration consists of finding the spatial transformation $T$ that best aligns
a \emph{moving} image $M$ with a \emph{fixed} image $F$. In this work, the fixed
image is the template and each AIIB23 segmentation is, in turn, the moving
image. The transformation is restricted to the class of three-dimensional affine
maps,
\begin{equation}
  T_{\boldsymbol{\mu}}(\mathbf{x}) = \mathbf{A}\,(\mathbf{x} - \mathbf{c}) + \mathbf{c} + \mathbf{t},
  \label{eq:affine}              % affine map: linear part A plus translation t
\end{equation}
where $\mathbf{x} \in \mathbb{R}^3$ is a point in the fixed image domain,
$\mathbf{A} \in \mathbb{R}^{3\times 3}$ is a general matrix,
$\mathbf{t} \in \mathbb{R}^3$ is a translation vector and $\mathbf{c}$ is a
fixed centre of rotation. The parameter vector $\boldsymbol{\mu}$ therefore
contains $9 + 3 = 12$ degrees of freedom, which jointly encode translation,
rotation, anisotropic scaling and shear. Because the transformation is linear up
to a translation, straight lines are mapped to straight lines and parallelism is
preserved, so the shape of the airway tree is modified only globally.

% ------------------------------------------------------------
\paragraph{Similarity metric and optimisation.}  % what is optimised and how

The optimal parameters are obtained by solving
\begin{equation}
  \hat{\boldsymbol{\mu}} = \arg\min_{\boldsymbol{\mu}} \;
  \mathcal{C}\big(F,\, M \circ T_{\boldsymbol{\mu}}\big),
  \label{eq:optimisation}        % registration as an optimisation problem
\end{equation}
where $\mathcal{C}$ is a cost function measuring the dissimilarity between the
fixed image and the transformed moving image. The cost function used is the
negative of the Mattes mutual information \cite{mattes2003petct},
\begin{equation}
  \mathcal{C}(\boldsymbol{\mu}) = -\sum_{\iota}\sum_{\kappa} p(\iota,\kappa;\boldsymbol{\mu})
  \log \frac{p(\iota,\kappa;\boldsymbol{\mu})}{p_F(\iota)\, p_M(\kappa;\boldsymbol{\mu})},
  \label{eq:mattes}              % negative mutual information from the joint histogram
\end{equation}
where $p(\iota,\kappa;\boldsymbol{\mu})$ is the joint probability of intensity
bins $\iota$ and $\kappa$ in the fixed and transformed moving images, and $p_F$
and $p_M$ are the corresponding marginal distributions. In the Mattes
formulation, these distributions are estimated with B-spline Parzen windows,
which makes the metric differentiable with respect to $\boldsymbol{\mu}$. The
histograms were built with 32 bins from a random sample of 20\% of the voxels,
which reduces the computational cost of each iteration.

Equation~\eqref{eq:optimisation} was solved with a gradient descent optimiser.
Before optimisation, the transformation was initialised by aligning the
intensity centres of mass of the two images, which provides a coarse
translational alignment and reduces the risk of convergence to a poor local
optimum.

% ------------------------------------------------------------
\paragraph{Multi-resolution strategy.}  % coarse-to-fine scheme

To improve robustness and speed, the optimisation was performed within a
four-level multi-resolution scheme \cite{thevenaz1998pyramid}. At each level,
both images were smoothed with a Gaussian kernel and downsampled by an integer
shrink factor; the optimisation then proceeded from the coarsest level to the
original resolution, with the result of each level used to initialise the next.
Coarse levels capture the large-scale alignment, while finer levels refine it.
The settings of each level are listed in Table~\ref{tab:multires}.

\begin{table}[htbp]              % per-level settings of the multi-resolution pyramid
  \centering                     % centre the table on the page
  \caption{Multi-resolution settings of the affine registration.}
  \label{tab:multires}           % referenced in the text above
  \begin{tabular}{cccc}          % four centred columns
    \toprule
    Level & Shrink factor & Gaussian $\sigma$ (voxels) & Maximum iterations \\
    \midrule
    1 (coarsest) & 6 & 3 & 2100 \\
    2            & 4 & 2 & 1200 \\
    3            & 2 & 1 & 1200 \\
    4 (finest)   & 1 & 0 & 10   \\
    \bottomrule
  \end{tabular}
\end{table}

At each level, the optimisation stopped either when the maximum number of
iterations was reached or when the cost function had converged, i.e.\ when its
variation over the last 10 iterations fell below $10^{-6}$.

% ------------------------------------------------------------
\paragraph{Resampling of the registered images.}  % how outputs were produced

Once the optimal transformation $T_{\hat{\boldsymbol{\mu}}}$ had been estimated,
each moving segmentation was resampled onto the voxel grid of the template using
linear interpolation and saved in NIfTI format. Since the inputs are binary
masks, linear interpolation yields intermediate values in $[0,1]$ at the
boundary of the airway lumen, while voxels well inside or outside the airway
keep the values 1 and 0 respectively.
% TODO: if the outputs were thresholded afterwards (e.g. at 0.5), state it here.

% ------------------------------------------------------------
\paragraph{Implementation.}      % software and computing environment

The registration was implemented in Python~3.11 with ANTsPy (version
\antspyversion) \cite{tustison2021antsx}, the Python interface of the Advanced
Normalization Tools (ANTs) \cite{avants2009ants}, which is built on the
registration framework of the Insight Toolkit (ITK) \cite{avants2014itkv4}. The
affine registration was obtained with the function \texttt{ants.registration}
using the predefined \texttt{Affine} transformation type, whose default settings
correspond to the parameters described above under \emph{Similarity metric and
optimisation} and \emph{Multi-resolution strategy}.

Since each segmentation is registered to the template independently of the
others, the task is embarrassingly parallel. The computation was run on a single
node of a high-performance computing cluster managed by the SLURM workload
manager \cite{yoo2003slurm}, with 48 CPU cores and 128\,GB of memory allocated
to the job. Parallelism was organised on two levels: the joblib library
\cite{joblib} distributed the registrations over 12 independent worker
processes, and each registration used 4 ITK threads, so that
\begin{equation}
  12 \text{ processes} \times 4 \text{ threads} = 48 \text{ cores},
  \label{eq:cores}               % total cores used = processes x threads per process
\end{equation}
matching the allocated resources exactly. Segmentations for which a registered
output already existed were skipped, which allowed an interrupted job to be
resumed without repeating completed registrations.

Because the similarity metric is estimated from a random subset of voxels and no
fixed random seed was set, repeating a registration may yield slightly different
transformation parameters; these differences are expected to be small relative
to the global misalignment corrected by the affine transformation.
